% TOP_PROBLEM: {"id": 2306077, "title": "Research Problems in Function Theory — Problem 6.77", "category_id": 8, "category": {"id": 8, "name": "analysis", "display_name": "Analysis", "description": "Limits, continuity, calculus, and function theory.", "slug": "analysis", "order_index": 8, "created_at": "2026-07-31T15:26:25.671Z"}, "status": "open"}
% FIRST_POSTED: null
% ENTRY_KIND: statement_only
% Imported from: batch03.tex; UnsolvedMath ID 2306077
\documentclass[11pt]{article}
\usepackage[margin=25mm]{geometry}
\usepackage{fontspec}
\setmainfont{Latin Modern Roman}
\usepackage{amsmath,amssymb,mathtools}
\usepackage{unicode-math}
\setmathfont{Latin Modern Math}
\usepackage{xurl,hyperref}
\hypersetup{colorlinks=true,urlcolor=blue}
\setcounter{secnumdepth}{0}
\setlength{\parindent}{0pt}
\setlength{\parskip}{5pt}
\setlength{\emergencystretch}{3em}
\newcommand{\sref}[2]{\hyperlink{source:#1:#2}{[S#2]}}
\newcommand{\cataloguescope}[1]{\paragraph{Recorded target.}#1}
\begin{document}
% UnsolvedMath ID 2306077; source catalogue code AMR-022-6077
\setcounter{section}{2306076}
\section{Research Problems in Function Theory — Problem 6.77}\label{2306077}
{\small\sffamily UnsolvedMath ID 2306077 · Analysis}
\subsection{Definitions and mathematical statement}

\paragraph{Shared notation.}$\mathbb N=\{1,2,\ldots\}$, and $\mathbb Z,\mathbb Q,\mathbb R,\mathbb C$ denote the integers, rationals, reals and complex numbers. Any other conventions are specified locally.


Write $\mathbb D=\{z\in\mathbb C:|z|<1\}$ and $\mathbb T=\partial\mathbb D$. For $n\in\mathbb N$, let $\mathcal P_n$ consist of polynomials $p(z)=z+\sum_{j=2}^n a_jz^j$ of degree at most $n$, injective on $\mathbb D$. Their boundary values are ordinary polynomial values on $\mathbb T$, and integration is with respect to angular Lebesgue measure.
\paragraph{Statement recorded in the supplied source.}
For every $n\in\mathbb N$ and $0<q<\infty$, determine \[\max_{p\in\mathcal P_n}\int_0^{2\pi}|p(e^{it})|^q\,dt.\] \sref{2306077}{2}
\subsection{Short English statement}
Find the largest boundary integral mean of every positive exponent among normalized univalent polynomials of bounded degree.
\subsection{Sources}
\begin{description}
\item[\hypertarget{source:2306077:1}{S1}] ULAM.AI and the UnsolvedMath Contributors, \emph{UnsolvedMath: A Curated Collection of Open Mathematics Problems}, record 2306077 (AMR-022-6077). CC BY 4.0. \url{https://huggingface.co/datasets/ulamai/UnsolvedMath}.
\item[\hypertarget{source:2306077:2}{S2}] W. K. Hayman and E. F. Lingham, \emph{Research Problems in Function Theory} (2018), Chapter 6, Problem 6.77 and its complete update; Chapter 6 initial notation. \url{https://arxiv.org/abs/1809.07200}.
\end{description}
\label{end:2306077}
\end{document}
